Cartesian coordinates specify positions directly in space, usually in angstroms (\AA) in the modeling interface. Fractional coordinates give proportions along the three lattice vectors. For $(u,v,w)$, the position is $u\mathbf a+v\mathbf b+w\mathbf c$. These components are dimensionless: they specify how far to go along each lattice vector, rather than distances along global Cartesian axes.

To read the lattice input table, arrange the vectors as matrix rows:
\[
 A=\begin{bmatrix}
 a_x&a_y&a_z\\
 b_x&b_y&b_z\\
 c_x&c_y&c_z
 \end{bmatrix},\qquad
 \mathbf f=(u,v,w),\qquad \mathbf r=(x,y,z).
\]
The rows are $\mathbf a$, $\mathbf b$, and $\mathbf c$; the columns within each row are global $x,y,z$ components. With the row-vector convention used in this manual,
\begin{equation}
 \mathbf r=\mathbf f A,\qquad \mathbf f=\mathbf r A^{-1}.
 \label{eq:row-lattice-coordinates}
\end{equation}
Here $A^{-1}$ is the inverse of the complete lattice matrix, not the reciprocal of each entry. The cell must have nonzero volume for a Cartesian position to determine unique fractional coordinates.

For a simple example, take $\mathbf a=(4,0,0)$, $\mathbf b=(1,3,0)$, and $\mathbf c=(0,0,6)$, all in \AA. The fractional position $(0.5,0.5,0)$ gives
\[
 \mathbf r=0.5(4,0,0)+0.5(1,3,0)
          =(2.5,1.5,0)\,\text{\AA}.
\]
Half of $\mathbf b$ includes $0.5\,\text{\AA}$ along global $x$ and $1.5\,\text{\AA}$ along global $y$. The second fractional component of a skewed cell is therefore not its Cartesian $y$ coordinate. Figure~\ref{fig:round10-fractional-cartesian} shows the addition as two joined arrows.

\begin{figure}[!htb]
\centering
\begin{tikzpicture}[x=1mm,y=1mm,font=\sffamily\small]
 \fill[paleblue] (8,6)--(56,6)--(68,42)--(20,42)--cycle;
 \draw[teal,line width=.7pt] (8,6)--(56,6)--(68,42)--(20,42)--cycle;
 \draw[->,>=stealth,muted] (8,6)--(84,6) node[right] {$x$};
 \draw[->,>=stealth,muted] (8,6)--(8,49) node[above] {$y$};
 \draw[->,>=stealth,brand,line width=1pt] (8,6)--(32,6);
 \draw[->,>=stealth,brand,line width=1pt] (32,6)--(38,24);
 \draw[->,>=stealth,teal,line width=.8pt] (8,6)--(38,24);
 \draw[dashed,muted] (38,24)--(38,6);
 \fill[brand] (38,24) circle (1.2);
 \node[anchor=north,text=brand] at (22,5) {$u\mathbf a$};
 \node[anchor=west,text=brand] at (38,14) {$v\mathbf b$};
 \node[anchor=south east,text=teal] at (25,16) {$\mathbf r$};
 \node[anchor=north east] at (7,5) {$O$};
 \node[anchor=west,text=brand] at (42,27) {$\mathbf f=(u,v,0)$};
 \node[anchor=west,text=muted,font=\sffamily\footnotesize] at (78,35) {Cell cross-section};
 \node[anchor=west,text=muted,font=\sffamily\footnotesize] at (78,28) {Add vectors for position};
 \node[anchor=north,text=teal] at (54,4) {$\mathbf a$};
 \node[anchor=east,text=teal] at (15,31) {$\mathbf b$};
\end{tikzpicture}
\caption{Fractional and Cartesian coordinates (2D schematic). Adding $u\mathbf a$ and $v\mathbf b$ gives $\mathbf r$. In a skewed cell, $\mathbf b$ has both global $x$ and $y$ components.}
\label{fig:round10-fractional-cartesian}
\end{figure}


In a cubic cell of side $5\,\text{\AA}$, $(0.5,0,0)$ instead gives a displacement of $2.5\,\text{\AA}$ along $\mathbf a$. The lattice differs, not the conversion rule. Holding $\mathbf f$ fixed while changing $A$ generally moves atoms with the cell. To keep their spatial positions fixed, retain $\mathbf r$ and recalculate $\mathbf f$ using the new lattice. Section~\ref{sec:lattice-transform-reference} in Chapter~\ref{ch:modeling-tools} explains these options.

\begin{table}[!htb]
\centering\small
\begin{tabularx}{\linewidth}{llX}
\toprule
Symbol & Units & Meaning in the parameters\\
\midrule
$A$ & \AA & Lattice matrix with rows $\mathbf a,\mathbf b,\mathbf c$.\\
$\mathbf f$ & Dimensionless & Fractional coordinates along the lattice vectors.\\
$\mathbf r,\mathbf p$ & \AA & Cartesian position and a rotation anchor or point on a mirror plane.\\
$\mathbf n,\hat{\mathbf n}$ & Direction & A nonzero direction vector and its unit vector; direction alone is not a length in \AA.\\
$\theta$ & Degrees & Dialog rotation angle, positive by the right-hand rule.\\
$S$ & Integers & Supercell transformation matrix defining the new lattice.\\
$N,V$ & Count, \AA$^3$ & Independent model-site count and cell volume.\\
$\varepsilon_a,\varepsilon_b,\varepsilon_c$ & Dimensionless & Relative extension along the three lattice vectors.\\
\bottomrule
\end{tabularx}
\caption{Symbols used in modeling formulas. Coordinates and anchors have length units; direction vectors specify orientation. Their roles in dialogs differ.}
\label{tab:round10-modeling-symbols}
\end{table}
